\section{Detailed derivations and implementation formulas}
\label{app:derivations}

\subsection{Local gradient update in AESP normalization}

For the coordinate step
\(
 x^+=x-2\omega\nabla_u f(x)e_u/(1+\alpha)
\),
\begin{align}
  \nabla f(x^+)
  &=\nabla f(x)+Q(x^+-x)
  \notag\\
  &=\nabla f(x)-\omega\nabla_u f(x)e_u
    +\chi\omega\nabla_u f(x)D^{-1/2}AD^{-1/2}e_u,
  \qquad
  \chi=\frac{1-\alpha}{1+\alpha}.
  \label{eq:app-gradient-update}
\end{align}
Multiplying by \(D^{1/2}\) and using
\(
D^{1/2}(D^{-1/2}AD^{-1/2})e_u=Ae_u/\sqrt{d_u}
\)
and \(\nabla_u f(x)=q_{\mathrm{sc},u}/\sqrt{d_u}\) yields
\begin{equation}
    q_{\mathrm{sc}}^+
    =q_{\mathrm{sc}}-\omega q_{\mathrm{sc},u}e_u
     +\chi\omega q_{\mathrm{sc},u}AD^{-1}e_u,
    \label{eq:app-scaled-update}
\end{equation}
Thus a sparse implementation changes only \(u\) and its neighbors:
\begin{equation}
    q_{\mathrm{sc},u}^+=(1-\omega)q_{\mathrm{sc},u},
    \qquad
    q_{\mathrm{sc},v}^+
    =q_{\mathrm{sc},v}+\chi\omega\frac{q_{\mathrm{sc},u}}{d_u}
    \quad(v\sim u).
    \label{eq:app-coordinate-updates}
\end{equation}

\subsection{Transfer to the Work-I residual}

At an AESP checkpoint \(x\), compute
\begin{equation}
    \widehat r
    =-\frac{2}{1+\alpha}\nabla f(x).
    \label{eq:app-handoff-residual}
\end{equation}
If the code stores degree-scaled quantities,
\begin{equation}
    D^{1/2}\widehat r
    =-\frac{2}{1+\alpha}q_{\mathrm{sc}}(x).
    \label{eq:app-handoff-scaled}
\end{equation}
The Work-I active condition
\(
\abs{\widehat r_u}\geq 2\alpha\epsppr\sqrt{d_u}/(1+\alpha)
\)
is then exactly \cref{eq:final-active-unscaled}.

\subsection{Why objective convergence alone does not bound weighted mass}

Smoothness gives
\begin{equation}
    \norm{\nabla f(x)}_2^2
    \leq2\,[f(x)-f(x^0)].
    \label{eq:gradient-gap-upper}
\end{equation}
If \(\Omega_x:=\supp(\nabla f(x))\), then
\begin{align}
    M(x)
    &=\sum_{u\in\Omega_x}\sqrt{d_u}\abs{\nabla_u f(x)}
    \notag\\
    &\leq
    \sqrt{\vol(\Omega_x)}\norm{\nabla f(x)}_2
    \notag\\
    &\leq
    \sqrt{2\vol(\Omega_x)[f(x)-f(x^0)]}.
    \label{eq:mass-from-gap}
\end{align}
Thus objective acceleration implies weighted-mass compression only after
controlling the degree-volume of the gradient support.  This is another form
of the locality bottleneck.

\subsection{Marginal handoff rule from the mass tail bound}

Suppose the current best AESP checkpoint has mass \(M_{t-1}\), the next stage
costs \(W_t\), and it improves the mass by
\begin{equation}
    \Delta M_t:=M_{t-1}-M_t>0.
    \label{eq:mass-improvement}
\end{equation}
The proved \(\omega=1\) tail upper bound decreases by
\begin{equation}
    \Delta \widehat W_{\rm tail}
    =\frac{1+\alpha}{2\alpha^2\epsppr}\Delta M_t.
    \label{eq:predicted-work-saved}
\end{equation}
A natural marginal rule is therefore to continue AESP only while
\begin{equation}
    W_t
    <
    \frac{1+\alpha}{2\alpha^2\epsppr}\Delta M_t.
    \label{eq:marginal-switch}
\end{equation}
This is a heuristic because the tail formula is an upper bound, not an exact
cost model.  It is nevertheless preferable to an arbitrary fixed number of
outer stages and directly reflects the empirically observed knee.

\subsection{Plain optimal SOR: convergence without a local mass theorem}

By \cref{lem:coordinate-decrease}, plain \(\omega_\star\)-SOR is an objective
descent method.  Its general tail bound is
\begin{equation}
    \Work_{\rm tail}^{\omega_\star}
    \leq
    \frac{(1+\alpha)[f(x_0)-f(x^0)]}
         {\omega_\star(2-\omega_\star)\alpha^2\epsppr^2}.
    \label{eq:omega-star-objective-work}
\end{equation}
The product is
\begin{equation}
    \omega_\star(2-\omega_\star)
    =\frac{8\sqrt\alpha(1+\alpha)}{(1+\sqrt\alpha)^4}.
    \label{eq:omega-star-product}
\end{equation}
This objective-decrease analysis is insufficient to establish the desired
local \(1/\epsppr\) dependence.  It proves convergence but does not capture the
multi-step spectral acceleration that makes optimal SOR fast in practice.

\subsection{AESP inner condition number as an effective damping parameter}

Write
\begin{equation}
    \nu:=\frac{\alpha+\kappa_{\mathrm A}}{1+\kappa_{\mathrm A}}.
    \label{eq:effective-damping}
\end{equation}
Then
\begin{equation}
    Q+\kappa_{\mathrm A} I
    =(1+\kappa_{\mathrm A})
      \left[
        \frac{1+\nu}{2}I-
        \frac{1-\nu}{2}D^{-1/2}AD^{-1/2}
      \right].
    \label{eq:shifted-ppr-form}
\end{equation}
For \(\kappa_{\mathrm A}=1-2\alpha\), \(\nu=1/2\).  Hence the inner system is
not merely constant-conditioned abstractly; it is exactly a uniformly damped
PageRank-type system.  This structural observation supports reusing local
PPR-style coordinate solvers inside AESP.
